\chapter{Nucleophilic Substitution}\label{ch:b1:nucleophilic-substitution}

Take a bottle of \cip{S}-2-bromobutane, a single enantiomer that rotates
polarised light. Heated with sodium hydroxide in a polar \omterm{def:b1:intermolecular-forces-solvents:solvent-classes}{aprotic solvent},
it gives butan-2-ol that is a single enantiomer too --- but the other
one, \cip{R}: the reaction has turned the molecule inside out like an
umbrella in the wind. Left instead in warm water, the same bromide gives
butan-2-ol that is partly racemic. The same overall change, an OH group
replacing a Br atom, follows two different paths. This chapter
establishes the two mechanisms from their \omterm{def:b1:rate-laws:order}{rate laws}, explains their
stereochemistry, and gives the rules that decide between them --- the
first full use of the tools of the kinetics and electronic-effects
chapters.

\begin{recall}
\omterm{def:b1:rate-laws:order}{Rate laws}, \omterm{def:b1:elementary-steps:elementary-step}{elementary steps}, the \omterm{def:b1:elementary-steps:rds}{rate-determining step} and the
\omterm{def:b1:elementary-steps:qssa}{steady-state approximation} are in \cref{ch:b1:rate-laws,ch:b1:elementary-steps};
\omterm{def:b1:stereochemistry-in-depth:isomers}{configurations}, Cram drawings and \omterm{def:b1:stereochemistry-in-depth:optical-activity}{optical activity} in
\cref{ch:b1:stereochemistry-in-depth}; \omterm{def:b1:electronic-effects:nucleophile}{nucleophiles}, \omterm{def:b1:electronic-effects:nucleophile}{electrophiles},
\omterm{def:b1:electronic-effects:nucleophile}{leaving groups} and \omterm{def:b1:electronic-effects:intermediates}{carbocations} in \cref{ch:b1:electronic-effects};
protic and \omterm{def:b1:intermolecular-forces-solvents:solvent-classes}{aprotic solvents} in \cref{ch:b1:intermolecular-forces-solvents}.
\end{recall}

\section{The substitution and its partners}

\begin{definition}[Nucleophilic substitution, substrate]\label{def:b1:nucleophilic-substitution:substitution}
A \emph{nucleophilic substitution}\index{nucleophilic substitution} is
the replacement, on a saturated carbon, of a \omterm{def:b1:electronic-effects:nucleophile}{leaving group} Y by a
\omterm{def:b1:electronic-effects:nucleophile}{nucleophile} Nu: $\mathrm{Nu^-} + \mathrm{R{-}Y} \to \mathrm{R{-}Nu} +
\mathrm{Y^-}$. The compound R--Y is the \emph{substrate}\index{substrate};
the carbon bearing Y is electrophilic because the C--Y bond is polarised
towards Y.
\end{definition}

Halogenoalkanes are the typical \omterm{def:b1:nucleophilic-substitution:substitution}{substrates}, with \ce{Cl-}, \ce{Br-} or
\ce{I-} as \omterm{def:b1:electronic-effects:nucleophile}{leaving groups}. Hydroxide gives alcohols; alkoxides give
ethers; cyanide gives nitriles (one more carbon); ammonia and amines give
amines; iodide exchanges with the other halides. Water and alcohols, used
as solvents, can also act as \omterm{def:b1:electronic-effects:nucleophile}{nucleophiles}: the substitution is then a
\omterm{def:b1:nucleophilic-substitution:sn1}{solvolysis}.

\section{The SN2 mechanism}

When bromomethane reacts with hydroxide in water, doubling either
concentration doubles the rate: $v = k[\ce{CH3Br}][\ce{OH-}]$, a reaction of
order one in each reactant.

\begin{definition}[SN2 mechanism, Walden inversion]\label{def:b1:nucleophilic-substitution:sn2}
The \emph{SN2 mechanism}\index{SN2 mechanism} (substitution,
nucleophilic, bimolecular) is a single \omterm{def:b1:elementary-steps:elementary-step}{elementary step} in which the
\omterm{def:b1:electronic-effects:nucleophile}{nucleophile} attacks the carbon from the side opposite the \omterm{def:b1:electronic-effects:nucleophile}{leaving group}
while the \omterm{def:b1:electronic-effects:nucleophile}{leaving group} departs. The three other substituents of the
carbon flip over, like an umbrella turned inside out: the \emph{Walden
inversion}\index{Walden inversion}.
\end{definition}

\begin{omfigure}
\setchemfig{atom sep=2.1em}
\schemestart
\chemfig{H@{nu}O^{-}}
\arrow{0}[,0.35]
\chemfig{@{c}C(-[:150]H_3C)(<[:210]H)(<:[:250]C_2H_5)-[@{cb}]@{br}Br}
\arrow{->}
\chemfig{[:0]HO-C(-[:30]CH_3)(<[:-30]H)(<:[:-70]C_2H_5)}
\+
\chemfig{Br^{-}}
\schemestop
\chemmove{\draw[->, thick, shorten <=2pt, shorten >=3pt](nu) .. controls +(15:5mm) and +(175:5mm) .. (c);
          \draw[->, thick, shorten <=2pt](cb) .. controls +(90:6mm) and +(90:6mm) .. (br);}

\medskip
\begin{tikzpicture}[font=\small]
  \node[draw, dashed, rounded corners, inner sep=6pt] at (0,0)
    {$\left[\vcenter{\hbox{\ce{HO}$\cdots$\chemfig{C(-[2]CH_3)(-[:-30]H)(-[:210]C_2H_5)}$\cdots$\ce{Br}}}\right]^{-}$};
  \node[below] at (0,-1.1) {transition state: five groups around C, three in a plane};
\end{tikzpicture}

{\small The SN2 reaction of hydroxide with \cip{S}-2-bromobutane. The
\omterm{def:b1:electronic-effects:nucleophile}{nucleophile} attacks the face opposite the bromine (\omterm{def:b1:electronic-effects:curly-arrow}{curly arrows}); in the
\omterm{def:b1:elementary-steps:energy-profile}{transition state} the O--C bond is forming while the C--Br bond breaks;
the product is \cip{R}-butan-2-ol: the three other groups have flipped to
the other side.}
\end{omfigure}

\begin{proposition}[Rate law of SN2]\label{prop:b1:nucleophilic-substitution:sn2-law}
For an SN2 reaction, $v = k[\mathrm{RY}][\mathrm{Nu}]$.
\end{proposition}

\begin{proof}
The mechanism is a single bimolecular \omterm{def:b1:elementary-steps:elementary-step}{elementary step}; by the law of mass
action for \omterm{def:b1:elementary-steps:elementary-step}{elementary steps} (\cref{ch:b1:elementary-steps}) its rate is
proportional to the product of the concentrations of the two partners.
\end{proof}

\begin{proposition}[Stereochemistry of SN2]\label{prop:b1:nucleophilic-substitution:sn2-inversion}
An SN2 reaction at a stereogenic carbon gives a single stereoisomer, with
inverted arrangement of the three unchanged groups: a reaction of this
kind is stereospecific.
\end{proposition}

\begin{proof}
The \omterm{def:b1:electronic-effects:nucleophile}{nucleophile} can only reach the empty side of the carbon, opposite the
\omterm{def:b1:electronic-effects:nucleophile}{leaving group}; the three other groups pass through a plane in the
\omterm{def:b1:elementary-steps:energy-profile}{transition state} and end on the other side. The new bond points where
the old one did not: each molecule of \omterm{def:b1:nucleophilic-substitution:substitution}{substrate} gives the mirror
arrangement of its groups. The descriptor usually changes from \cip{S}
to \cip{R} or back, but not always, since it depends on the ranks of the
\omterm{def:b1:electronic-effects:nucleophile}{nucleophile} and of the \omterm{def:b1:electronic-effects:nucleophile}{leaving group}.
\end{proof}

\begin{definition}[Stereospecific reaction]\label{def:b1:nucleophilic-substitution:stereospecific}
A reaction is a \emph{stereospecific reaction}\index{stereospecific
reaction} when stereoisomeric \omterm{def:b1:nucleophilic-substitution:substitution}{substrates} give stereoisomerically
different products, each \omterm{def:b1:nucleophilic-substitution:substitution}{substrate} giving its own product.
\end{definition}

\section{The SN1 mechanism}

2-Bromo-2-methylpropane reacts with water at a rate that does not depend
on the concentration of the \omterm{def:b1:electronic-effects:nucleophile}{nucleophile}: $v = k[(\ce{CH3})_3\ce{CBr}]$. A
step that does not involve the \omterm{def:b1:electronic-effects:nucleophile}{nucleophile} must decide the rate.

\begin{definition}[SN1 mechanism, racemisation, solvolysis]\label{def:b1:nucleophilic-substitution:sn1}
The \emph{SN1 mechanism}\index{SN1 mechanism} (substitution,
nucleophilic, unimolecular) has two steps: the slow \omterm{def:b1:electronic-effects:cleavage}{heterolysis} of the
C--Y bond, giving a \omterm{def:b1:electronic-effects:intermediates}{carbocation}, then the fast addition of the
\omterm{def:b1:electronic-effects:nucleophile}{nucleophile} to the \omterm{def:b1:electronic-effects:intermediates}{carbocation}. A reaction that turns a single
enantiomer into a mixture of both is a \emph{racemisation}\index{racemisation};
a substitution in which the solvent is the \omterm{def:b1:electronic-effects:nucleophile}{nucleophile} is a
\emph{solvolysis}\index{solvolysis}.
\end{definition}

\begin{omfigure}
\setchemfig{atom sep=2.1em}
\schemestart
\chemfig{C(-[2]CH_3)(-[4]H_3C)(-[6]CH_3)-[@{cb2}]@{br2}Br}
\arrow{->[slow]}
\chemfig{C^{+}(-[2]CH_3)(-[:210]H_3C)(-[:-30]CH_3)}
\+
\chemfig{Br^{-}}
\schemestop

\smallskip
\schemestart
\chemfig{C^{+}(-[2]CH_3)(-[:210]H_3C)(-[:-30]CH_3)}
\arrow{->[fast][\ce{H2O}]}
\chemfig{C(-[2]CH_3)(-[4]H_3C)(-[6]CH_3)-OH_2^{+}}
\arrow{->[\ce{-H+}]}
\chemfig{C(-[2]CH_3)(-[4]H_3C)(-[6]CH_3)-OH}
\schemestop
\chemmove{\draw[->, thick, shorten <=2pt](cb2) .. controls +(90:6mm) and +(90:6mm) .. (br2);}

\medskip
\begin{tikzpicture}[font=\small]
  \fill[omProp!15, draw=omProp] (0,0.2) ellipse (0.2 and 0.85);
  \fill[omProp!15, draw=omProp] (0,-0.2) ellipse (0.2 and 0.85);
  \draw[thick] (0,0) -- (-1.1,-0.2) node[left] {\ce{CH3}};
  \draw[thick] (0,0) -- (0.95,-0.45) node[right] {\ce{C2H5}};
  \draw[thick] (0,0) -- (0.4,0.5) node[above right] {H};
  \node[fill=white, inner sep=1pt] at (0,0) {\ce{C+}};
  \draw[->, thick, omDef] (-0.9,1.6) node[left] {\ce{H2O}} -- (-0.15,0.95);
  \draw[->, thick, omDef] (0.9,-1.6) node[right] {\ce{H2O}} -- (0.15,-0.95);
  \node[align=center] at (4.6,0) {attack on either face:\\ both configurations,\\ racemic if nothing else\\ differs between the faces};
\end{tikzpicture}

{\small SN1 \omterm{def:b1:nucleophilic-substitution:sn1}{solvolysis} of 2-bromo-2-methylpropane in water: slow ionisation
to the tertiary \omterm{def:b1:electronic-effects:intermediates}{carbocation}, then capture by water and loss of a proton.
Below: a \omterm{def:b1:stereochemistry-in-depth:stereocentre}{chiral} secondary \omterm{def:b1:electronic-effects:intermediates}{carbocation} (from 2-bromobutane) is planar,
and water adds on either face with equal probability.}
\end{omfigure}

\begin{proposition}[Rate law of SN1]\label{prop:b1:nucleophilic-substitution:sn1-law}
For the \omterm{def:b1:nucleophilic-substitution:sn1}{SN1 mechanism}, if the \omterm{def:b1:electronic-effects:intermediates}{carbocation} is captured much faster than it
returns to the \omterm{def:b1:nucleophilic-substitution:substitution}{substrate}, $v = k_1[\mathrm{RY}]$, independent of the
\omterm{def:b1:electronic-effects:nucleophile}{nucleophile}.
\end{proposition}

\begin{proof}
Steps: $\mathrm{RY} \rightleftharpoons \mathrm{R^+} + \mathrm{Y^-}$ ($k_1$,
$k_{-1}$), $\mathrm{R^+} + \mathrm{Nu} \to \mathrm{RNu}$ ($k_2$). The
\omterm{def:b1:electronic-effects:intermediates}{carbocation} is a \omterm{def:b1:electronic-effects:intermediates}{reactive intermediate}: by the \omterm{def:b1:elementary-steps:qssa}{steady-state
approximation} (\cref{ch:b1:elementary-steps}), $k_1[\mathrm{RY}] =
k_{-1}[\mathrm{R^+}][\mathrm{Y^-}] + k_2[\mathrm{R^+}][\mathrm{Nu}]$, so
$v = k_2[\mathrm{R^+}][\mathrm{Nu}] = \dfrac{k_1k_2[\mathrm{RY}][\mathrm{Nu}]}
{k_{-1}[\mathrm{Y^-}] + k_2[\mathrm{Nu}]}$. When $k_2[\mathrm{Nu}] \gg
k_{-1}[\mathrm{Y^-}]$ this reduces to $k_1[\mathrm{RY}]$: the first step is
rate-determining.
\end{proof}

\begin{proposition}[Stereochemistry of SN1]\label{prop:b1:nucleophilic-substitution:sn1-stereo}
An SN1 reaction at a stereogenic carbon gives both \omterm{def:b1:stereochemistry-in-depth:isomers}{configurations}: the
reaction is not stereospecific, and leads to \omterm{def:b1:nucleophilic-substitution:sn1}{racemisation} when nothing
distinguishes the two faces of the \omterm{def:b1:electronic-effects:intermediates}{carbocation}.
\end{proposition}

\begin{proof}
The \omterm{def:b1:electronic-effects:intermediates}{carbocation} is planar (\cref{ch:b1:electronic-effects}): its two
faces are mirror images. If the \omterm{def:b1:electronic-effects:nucleophile}{leaving group} has gone far before the
capture, the \omterm{def:b1:electronic-effects:nucleophile}{nucleophile} adds to either face at the same rate, giving the
two \omterm{def:b1:stereochemistry-in-depth:enantiomers}{enantiomers} in equal amounts. In practice the leaving ion still
shields one face for a while, and a slight excess of inversion is often
observed.
\end{proof}

\begin{omfigure}
\begin{tikzpicture}
\begin{axis}[
  width=0.8\linewidth, height=5.0cm, xmin=0, xmax=10, ymin=-30, ymax=110,
  axis x line=bottom, axis y line=left, xlabel={reaction coordinate},
  ylabel={energy}, xtick=\empty, ytick=\empty,
  label style={font=\small}, legend style={font=\scriptsize, draw=none, at={(0.98,0.98)}, anchor=north east}]
\addplot[omDef, very thick] table[x=x, y=E] {figdata/out/bachelor-1-energy-profiles-sn2.dat};
\addplot[omProp, very thick, dashed] table[x=x, y=E] {figdata/out/bachelor-1-energy-profiles-sn1.dat};
\legend{SN2: one step, SN1: two steps}
\node[font=\scriptsize, anchor=north] at (axis cs:5,68) {carbocation};
\node[font=\scriptsize, anchor=north] at (axis cs:0.6,-3) {RY};
\node[font=\scriptsize, anchor=south] at (axis cs:9.6,-18) {RNu};
\end{axis}
\end{tikzpicture}

{\small Schematic \omterm{def:b1:elementary-steps:energy-profile}{energy profiles}. SN2: a single \omterm{def:b1:elementary-steps:energy-profile}{transition state}, five
groups around carbon. SN1: a first, higher barrier (ionisation, the
\omterm{def:b1:elementary-steps:rds}{rate-determining step}), the \omterm{def:b1:electronic-effects:intermediates}{carbocation} in a shallow well, then a small
barrier to its capture.}
\end{omfigure}

\section{What decides the mechanism}

\begin{proposition}[Substrate]\label{prop:b1:nucleophilic-substitution:substrate}
Methyl and primary \omterm{def:b1:nucleophilic-substitution:substitution}{substrates} react by SN2; tertiary \omterm{def:b1:nucleophilic-substitution:substitution}{substrates} by SN1;
secondary \omterm{def:b1:nucleophilic-substitution:substitution}{substrates} by either, depending on the other factors. Allylic
and benzylic \omterm{def:b1:nucleophilic-substitution:substitution}{substrates} react fast by both.
\end{proposition}

\begin{proof}
SN2 needs room behind the carbon: each alkyl group in place of a hydrogen
crowds the \omterm{def:b1:elementary-steps:energy-profile}{transition state}, which carries five groups, and slows the
reaction; three groups \omterm{def:b1:periodicity:table}{block} it. SN1 needs a stable \omterm{def:b1:electronic-effects:intermediates}{carbocation}: the
order of stability (\cref{prop:b1:electronic-effects:carbocation-order})
makes ionisation fast for tertiary, possible for secondary, and too slow
for primary and methyl \omterm{def:b1:electronic-effects:intermediates}{carbocations}. Allylic and benzylic \omterm{def:b1:electronic-effects:intermediates}{carbocations}
are stabilised by resonance, and the same $\pi$ system stabilises the SN2
\omterm{def:b1:elementary-steps:energy-profile}{transition state}.
\end{proof}

\begin{proposition}[Leaving group]\label{prop:b1:nucleophilic-substitution:leaving-group}
Both mechanisms are faster the better the \omterm{def:b1:electronic-effects:nucleophile}{leaving group}, that is the
weaker the base it becomes: $\ce{I-} > \ce{Br-} > \ce{Cl-} \gg \ce{F-}$;
\ce{OH-} and \ce{RO-} do not leave.
\end{proposition}

\begin{proof}
In both, the C--Y bond is broken in or before the \omterm{def:b1:elementary-steps:rds}{rate-determining step},
and the \omterm{def:b1:elementary-steps:energy-profile}{transition state} carries part of the negative charge on Y. A
group that holds a negative charge well --- the conjugate base of a
\omterm{def:b1:predominant-reaction:strong-acid}{strong acid} --- lowers that \omterm{def:b1:elementary-steps:energy-profile}{transition state}. \ce{HI}, \ce{HBr} and
\ce{HCl} are \omterm{def:b1:predominant-reaction:strong-acid}{strong acids}, \ce{HF} a weak one ($\mathrm{p}K_a$ 3.2), water and
alcohols very weak ones: \ce{OH-} and \ce{RO-} are \omterm{def:b1:predominant-reaction:strong-acid}{strong bases} and poor
\omterm{def:b1:electronic-effects:nucleophile}{leaving groups}.
\end{proof}

\begin{proposition}[Solvent]\label{prop:b1:nucleophilic-substitution:solvent}
Polar \omterm{def:b1:intermolecular-forces-solvents:solvent-classes}{protic solvents} (water, alcohols) favour SN1; polar \omterm{def:b1:intermolecular-forces-solvents:solvent-classes}{aprotic solvents}
(propanone, dimethyl sulfoxide, dimethylformamide) favour SN2 with anionic
\omterm{def:b1:electronic-effects:nucleophile}{nucleophiles}.
\end{proposition}

\begin{proof}
SN1 creates two ions from a neutral \omterm{def:b1:nucleophilic-substitution:substitution}{substrate}: a polar \omterm{def:b1:intermolecular-forces-solvents:solvent-classes}{protic solvent}
stabilises both, the cation through its \omterm{def:b1:lewis-resonance-vsepr:lewis}{lone pairs} and the anion through
\omterm{def:b1:intermolecular-forces-solvents:hbond}{hydrogen bonds}, and lowers the ionisation barrier. In SN2 the charge of
an anionic \omterm{def:b1:electronic-effects:nucleophile}{nucleophile} is spread in the \omterm{def:b1:elementary-steps:energy-profile}{transition state}; a \omterm{def:b1:intermolecular-forces-solvents:solvent-classes}{protic
solvent}, which binds the small anion strongly by \omterm{def:b1:intermolecular-forces-solvents:hbond}{hydrogen bonds}, slows
it, whereas an \omterm{def:b1:intermolecular-forces-solvents:solvent-classes}{aprotic solvent} leaves the anion free and reactive
(\cref{ch:b1:intermolecular-forces-solvents}).
\end{proof}

\begin{method}[Choosing the mechanism]\label{met:b1:nucleophilic-substitution:choose-mechanism}
\begin{enumerate}
  \item \omterm{def:b1:nucleophilic-substitution:substitution}{Substrate}: methyl or primary $\to$ SN2; tertiary $\to$ SN1;
        secondary $\to$ go on.
  \item \omterm{def:b1:electronic-effects:nucleophile}{Nucleophile}: strong and anionic ($\ce{OH-}$, $\ce{RO-}$, $\ce{CN-}$,
        $\ce{I-}$) $\to$ SN2; weak and neutral (water, alcohol) $\to$ SN1.
  \item Solvent: aprotic $\to$ SN2; protic $\to$ SN1.
  \item Predict the stereochemistry (inversion or \omterm{def:b1:nucleophilic-substitution:sn1}{racemisation}) and check
        that a \omterm{def:b1:predominant-reaction:strong-acid}{strong base} does not rather cause an elimination
        (\cref{ch:b1:elimination}).
\end{enumerate}
\end{method}

\begin{omfigure}
\begin{tabular}{lll}
\hline
 & SN2 & SN1 \\
\hline
steps & one & two, via a \omterm{def:b1:electronic-effects:intermediates}{carbocation} \\
\omterm{def:b1:rate-laws:order}{rate law} & $k[\mathrm{RY}][\mathrm{Nu}]$ & $k[\mathrm{RY}]$ \\
\omterm{def:b1:nucleophilic-substitution:substitution}{substrate} & methyl $>$ primary $>$ secondary & tertiary $>$ secondary \\
\omterm{def:b1:electronic-effects:nucleophile}{nucleophile} & strong, anionic & weak, neutral (often the solvent) \\
solvent & polar aprotic & polar protic \\
stereochemistry & inversion (stereospecific) & \omterm{def:b1:nucleophilic-substitution:sn1}{racemisation} (mostly) \\
\hline
\end{tabular}

\smallskip
{\small The two mechanisms side by side.}
\end{omfigure}

\section{Exercises}

\begin{exercise}[$\star$]\label{exo:b1:nucleophilic-substitution:1}
Write the products of: bromoethane with sodium hydroxide; iodomethane with
sodium ethoxide; 1-chloropropane with potassium cyanide; bromomethane with
ammonia (first step). Identify \omterm{def:b1:nucleophilic-substitution:substitution}{substrate}, \omterm{def:b1:electronic-effects:nucleophile}{nucleophile} and \omterm{def:b1:electronic-effects:nucleophile}{leaving group}.
\end{exercise}

\begin{exercise}[$\star$]\label{exo:b1:nucleophilic-substitution:2}
For a reaction $\mathrm{RY} + \mathrm{Nu}$ with $v = k[\mathrm{RY}][\mathrm{Nu}]$,
what happens to the \omterm{def:b1:rate-laws:initial-rate}{initial rate} if $[\mathrm{Nu}]$ is tripled? If both
concentrations are halved? Same questions if $v = k[\mathrm{RY}]$.
\end{exercise}

\begin{exercise}[$\star$]\label{exo:b1:nucleophilic-substitution:3}
Predict the mechanism (SN1 or SN2): 1-bromobutane with sodium iodide in
propanone; 2-bromo-2-methylpropane in water; bromomethane with hydroxide;
2-bromobutane in ethanol without added base.
\end{exercise}

\begin{exercise}[$\star$]\label{exo:b1:nucleophilic-substitution:4}
Draw the product of the SN2 reaction of cyanide with \cip{R}-2-bromobutane
in \omterm{def:b1:stereochemistry-in-depth:representations}{Cram representation} and give its descriptor.
\end{exercise}

\begin{exercise}[$\star\star$]\label{exo:b1:nucleophilic-substitution:5}
Write the full \omterm{def:b1:nucleophilic-substitution:sn1}{SN1 mechanism} of the hydrolysis of 2-bromo-2-methylpropane,
with \omterm{def:b1:electronic-effects:curly-arrow}{curly arrows}, including the loss of the proton. Why is the \omterm{def:b1:rate-laws:order}{rate law}
first order although three species take part?
\end{exercise}

\begin{exercise}[$\star\star$]\label{exo:b1:nucleophilic-substitution:6}
Explain why 1-bromo-2,2-dimethylpropane, a primary bromide, reacts very
slowly by SN2 and does not react by SN1.
\end{exercise}

\begin{exercise}[$\star\star$]\label{exo:b1:nucleophilic-substitution:7}
\cip{S}-2-iodooctane loses its \omterm{def:b1:stereochemistry-in-depth:optical-activity}{optical activity} when kept in propanone
with sodium iodide, although no new compound forms. Explain, and say why
the rate of loss of activity is twice the rate of substitution.
\end{exercise}

\begin{exercise}[$\star\star$]\label{exo:b1:nucleophilic-substitution:8}
Rank by rate of SN2 with iodide: 2-bromo-2-methylpropane, bromomethane,
2-bromopropane, 1-bromopropane. Rank the same compounds by rate of SN1
\omterm{def:b1:nucleophilic-substitution:sn1}{solvolysis}.
\end{exercise}

\begin{exercise}[$\star\star$]\label{exo:b1:nucleophilic-substitution:9}
An optically pure secondary bromide undergoes \omterm{def:b1:nucleophilic-substitution:sn1}{solvolysis}; the product
rotates light at \qty{12}{\percent} of the value expected for pure
inversion. Give the proportions of the two \omterm{def:b1:stereochemistry-in-depth:enantiomers}{enantiomers} of the product and
interpret.
\end{exercise}

\begin{exercise}[$\star\star\star$]\label{exo:b1:nucleophilic-substitution:10}
Derive the general SN1 \omterm{def:b1:rate-laws:order}{rate law} with the \omterm{def:b1:elementary-steps:qssa}{steady-state approximation}, and
show that adding the \omterm{def:b1:electronic-effects:nucleophile}{leaving group} ion \ce{Y-} at the start slows the
reaction (the \omterm{def:b1:precipitation:common-ion}{common-ion effect} of kinetics). Under which condition is
this effect negligible?
\end{exercise}

\begin{exercise}[$\star\star\star$]\label{exo:b1:nucleophilic-substitution:11}
Two half-lives for the reaction of bromomethane with hydroxide: with
$[\ce{OH-}]_0 = [\ce{CH3Br}]_0 = C_0$, show that $1/[\ce{CH3Br}] = 1/C_0 +
kt$ and give $t_{1/2}$. With $[\ce{OH-}]_0 = 50\,C_0$, show that the
decay is exponential and give $t_{1/2}$. Data of the exercise: $k =
\qty{2.0e-4}{L/(mol.s)}$, $C_0 = \qty{0.010}{mol/L}$.
\end{exercise}

\begin{exercise}[$\star\star\star$]\label{exo:b1:nucleophilic-substitution:12}
Explain why an alcohol, \ce{ROH}, does not react with sodium bromide, but
reacts with concentrated hydrobromic acid to give \ce{RBr}. Write the
mechanism for a tertiary and for a primary alcohol.
\end{exercise}

\section{Problem: Solvolysis of tert-Butyl Chloride}

\begin{problem}[{Weekend problem --- order and rate constant from
conductivity data, the half-life, the mechanism and its stereochemistry,
and the activation energy of the solvolysis}]
\label{pb:b1:nucleophilic-substitution:1}
2-Chloro-2-methylpropane (tert-butyl chloride) is dissolved, at a small
concentration, in a water--propanone mixture. Its \omterm{def:b1:nucleophilic-substitution:sn1}{solvolysis}
\ce{(CH3)3CCl + H2O -> (CH3)3COH + H+ + Cl-} produces ions, and the
conductivity $\kappa$ of the solution, which is proportional to the
amount of ions formed, is recorded. After a very long time it reaches
$\kappa_\infty = \qty{2.000}{mS/cm}$ at both temperatures. Measurements
(\unit{mS/cm}):

\begin{center}
\begin{tabular}{lcccccccc}
\hline
$t$ (min) & 0 & 5 & 10 & 15 & 20 & 30 & 45 & 60 \\
\hline
$\kappa$, \qty{25.0}{\celsius} & 0 & 0.172 & 0.329 & 0.473 & 0.605 & 0.835 & 1.110 & 1.321 \\
$\kappa$, \qty{35.0}{\celsius} & 0 & 0.507 & 0.885 & 1.168 & 1.379 & 1.654 & 1.856 & 1.940 \\
\hline
\end{tabular}
\end{center}
% data generated in figdata/bachelor-1/solvolysis.py (story data)

\medskip
\noindent\textbf{Part I --- The order.}
\begin{enumerate}
  \item Why is the conductivity proportional to the amount of \omterm{def:b1:nucleophilic-substitution:substitution}{substrate}
        that has reacted?
  \item Express $[\mathrm{RCl}]/[\mathrm{RCl}]_0$ in terms of $\kappa$ and
        $\kappa_\infty$.
  \item Write the integrated law of a first-order reaction and the
        quantity to plot against $t$ to test it.
  \item Compute $\ln(\kappa_\infty - \kappa)$ at \qty{25.0}{\celsius} for each
        time.
  \item Is the plot a straight line? Conclude on the order.
  \item Water is in large excess. What does that hide in the order?
\end{enumerate}

\medskip
\noindent\textbf{Part II --- \omterm{def:b1:rate-laws:order}{Rate constant} and \omterm{def:b1:rate-laws:half-life}{half-life}.}
\begin{enumerate}[resume]
  \item Find the \omterm{def:b1:rate-laws:order}{rate constant} at \qty{25.0}{\celsius}, in \unit{s^{-1}}.
  \item Compute the \omterm{def:b1:rate-laws:half-life}{half-life}.
  \item At what time is \qty{90}{\percent} of the \omterm{def:b1:nucleophilic-substitution:substitution}{substrate} consumed?
  \item Find the \omterm{def:b1:rate-laws:order}{rate constant} at \qty{35.0}{\celsius}.
  \item By what factor has the rate increased for ten degrees?
  \item Check one value of the \qty{35}{\celsius} series against the law.
\end{enumerate}

\medskip
\noindent\textbf{Part III --- The mechanism.}
\begin{enumerate}[resume]
  \item Which mechanism do the \omterm{def:b1:nucleophilic-substitution:substitution}{substrate} and the solvent suggest?
  \item Write it with \omterm{def:b1:electronic-effects:curly-arrow}{curly arrows}.
  \item Show that it agrees with the order found.
  \item Adding sodium hydroxide at \qty{0.01}{mol/L} hardly changes the
        \omterm{def:b1:rate-laws:rate}{rate of disappearance} of the chloride. Why?
  \item The same experiment with the \omterm{def:b1:stereochemistry-in-depth:stereocentre}{chiral} tertiary chloride
        3-chloro-3-methylhexane, optically pure, gives an alcohol of almost
        zero \omterm{def:b1:stereochemistry-in-depth:optical-activity}{optical activity}. Explain.
  \item Draw an \omterm{def:b1:elementary-steps:energy-profile}{energy profile} of the reaction and mark the
        \omterm{def:b1:elementary-steps:rds}{rate-determining step}.
  \item Why would 1-chlorobutane hardly react in the same conditions?
\end{enumerate}

\medskip
\noindent\textbf{Part IV --- \omterm{def:b1:rate-laws:activation-energy}{Activation energy}.}
\begin{enumerate}[resume]
  \item Write the \omterm{thm:b1:rate-laws:arrhenius}{Arrhenius law}.
  \item Express $E_a$ from the two \omterm{def:b1:rate-laws:order}{rate constants}.
  \item Compute $E_a$ ($R = \qty{8.314}{J/(K.mol)}$).
  \item What does this energy measure, in the mechanism?
  \item Predict the \omterm{def:b1:rate-laws:order}{rate constant} at \qty{45.0}{\celsius}.
  \item State the \omterm{def:b1:rate-laws:activation-energy}{activation energy} of the \omterm{def:b1:nucleophilic-substitution:sn1}{solvolysis}, to two significant
        figures.
\end{enumerate}
\end{problem}
